\chapterNote{Based on a talk given on January 16, 2023}

\chapter{Our World Through Perception}
\tocChapterNote{Based on a talk given on January 16, 2023}

Yesterday was the anniversary of my mother's death six years ago, and I spent a little bit of time today looking at some photos. For the memorial service, my brother, sister, and I, along with various other family members, had put together a slideshow, a brief synopsis of her life in photos. I was struck once again at the variety of images, perceptions, and people that constitute what it is that I call my ``mother.''

There were photos of her as a young girl, of her growing up into young adulthood, and photos of her on the beach, probably in her twenties. She had dark brown hair and was a fairly attractive young woman. There were some wedding photos and some of my mother and father together in their early years. She started to have kids, and there were pictures of her feeding my brother as a baby. She was looking a little older and maybe a bit tired, but still lovingly feeding him. Then there were photos of her and the three of us, my brother, my sister, and I, as very young kids. She was getting into her thirties.

As she got older, into her middle and later years, the photos were of her with her grandchildren and my sister and some of myself in my thirties. The photos then progressed through her aging, final years, illness, and finally a number of photos of her after she passed away in her bed in the nursing home, where she spent the last six weeks of her life.

I found myself contemplating, ``Who was she? Who was this woman named Glenn DeMaria?'' (I was named after her, which caused a lot of confusion for people).

Who was this person that we label as ``Glenn DeMaria''? Was she that young girl on the horse, the mother with three children, the older woman in the wheelchair in her final days, or the lifeless corpse on a hospital bed? The worlds that we live in are worlds of perceptions, labels, and constructions of supposed realities. The people that inhabit these worlds are the same, in terms of experience. They are perceptions.

We create our perceptions of the world around us, and then we inhabit them. We take them all as so solid, so real. But right now, the entire lifetime of some person is \mbox{just a memory.}

\enlargethispage{1\baselineskip}
We can think about that in relation to ourselves as well. ``Who am I? Who was I?'' Looking at some of those photos, there was a young kid, maybe three or four years old, with a completely different body than this one. Blonde, straight hair in the early years, curly hair and a bushy beard later on, and then a body that morphed into middle age. Then there were some photos of me with my mother in her last few weeks, when I was a sixty-one-year-old monk. I think, ``Okay, well, who was that? Where is he now? What was that all about?'' It's just a memory now.

You can ask that of yourself. Think back to your early life. Go through the various stages. You have a completely different body now. Body cells replace themselves more or less every seven years, so you don't even have the same body that you had seven years ago. Who was Lee Mintz? I remember a Lee Mintz who had big, bushy, dark hair and was fifteen or twenty years younger than he is now. Who was he before that? Or Ty Krug? Or Nick Teal? Josh Dultz? Who were you? Where were you? I don't even remember Tan Sampanno's former name. I just remember a Korean monk who showed up at Abhayagiri. Where did they go? Just memories. I can remember Debbie Stamp from at least thirty years ago. Where did she go? These are questions for reflection that challenge our assumptions of the world we create around us. How real is it? How solid is it? We take it on as so real, so solid, but everything before this very moment is just a memory. The future is just a projection. Experience is constantly moving, changing, shifting.

Why do we take it all so seriously? For my mother, there are no more memories. Assuming she has been reborn, ``she'' \linebreak is experiencing another existence. Six years later, if she \linebreak has been reborn into the human realm, she is sitting somewhere as a six-year-old child. She probably has no memory \linebreak of Glenn DeMaria.

\looseness=-1
\enlargethispage{1.5\baselineskip}
Many Buddhist teachings instruct us to take a look at our perceptual world, the reality that we've assumed into existence. Those teachings encourage us to stop identifying with all of these constructs, to let the sense of ownership that comes with identification fade away. When we take on those constructs as who we are, that's the basis for suffering, \emph{dukkha}.

The Buddha's teachings point towards recognizing the arising of dukkha. How can we reshape our perceptions such that we can relax the process of grasping that is based on mistaken perceptions? Perceptions are based on memory. Perception, the aggregate ``\emph{saññā},'' is something with which we often strongly identify. ``Recognition'' might be a better translation, in that all of our perceptions, the ways that we experience the world, are through filters. Filters cause us to remember certain related experiences from the past that we label and categorize to create a sense of continuity. We need perceptions to navigate through the world without going insane because of lots of random sensory stimulation. We assign meaning to memories and experiences. ``Felt meaning'' is another translation of saññā that I have heard. Recognition is a cognizing process. It's based on memories that help us to create categories of people, situations, and mental structures. That's what we do to get through life without losing our sanity. But the perceptions we create are up to us, and so many times we create perceptual patterns that result in suffering.

Some of the teachings that the Buddha gives us are approaches to reconstruct perceptions in ways that are more suitable for living a happy, contented life, for letting go of dukkha. We've been recently having some readings related to the fourth \emph{satipaṭṭhāna}, the establishing of mindfulness of specific Dhamma themes. We've started the section on the five hindrances, discussing what they are and how we can work with them in our meditation and in our lives. The five hindrances are sensual desire, ill will, dullness/drowsiness, agitation/restlessness, and doubt. They are embedded habits that rise up when we're in contact with the world. They operate with craving, aversion, or confusion as their root causes and lead to dukkha. They're unsatisfactory, they're unpleasant, and there is the wish to be \mbox{free from them.}

These patterns drive our whole life, and we can choose to change them. The Buddha says that through appropriate attention, we can change our perceptual patterns and our experience of the world.

For instance, we don't have to experience the world through the five hindrances. We don't have to seek sensual pleasure in order to have a moment of contentment. We think that engaging in sensual pleasure is enjoyable, and to a certain extent it is, but the actual experience of sensual desire is something I find quite unpleasant. There is an agitating quality to wanting something.

When we compare it to a peaceful and contented mind, the arising of sensual desire is unpleasant. Even if it results in getting something in the sensual realm that's pleasurable, the process is agitating. The goal of sense desire is to get rid of the unpleasant feeling of wanting. And then when we do get some pleasure, it goes away very quickly. All in all, the pursuit of pleasure due to sensual desire is not gratifying in the long term.

All of the hindrances lead to dukkha. The hindrances are unpleasant ways of perceiving and living in the world, default patterns that we've learned and reinforced over a long period of time.

We can start to undo these patterns and reconstruct the \linebreak ways we perceive the world. We will also become more pleasant to be around. We all know what it's like to be around somebody who's consumed with craving, aversion, or other overt forms of the hindrances. It can be draining to be around them, just as it can be hard to be around ourselves when we're experiencing these hindrances.

Reconstructing these patterns is a gift to the world as \linebreak well as a gift to ourselves. We need to recognize them first as they arise in our practice and lives. We can't just pretend they're not there or try to paint over or annihilate them. It doesn't work. We have to first experience these patterns. That's the training in the four \emph{satipaṭṭhānas}: to experience the hindrances, know when they're present, know them very deeply and directly. It's not only a conceptual knowing but also allowing ourselves a direct experience of the hindrance. How does it feel in the body? How does it manifest in the mind, the heart? What are its qualities? Is it pleasant or is it painful? Does it cause the arising of suffering? To know experience directly is to know the First Noble Truth.

We have to know and understand the presence of dukkha. After we've seen and validated the experience of a hindrance, we can figure out how to apply an antidote to it. We can observe its arising and passing away. We can reframe our perceptions. We can look at an experience in a different way and turn it around. That's letting go of inappropriate attention and engaging in appropriate attention. For instance, we can use \emph{asubhakammaṭṭhāna} (contemplation of the unattractiveness of the body) as an antidote to sensual desire, and we can develop perceptions of goodwill and patience in response to the hindrance of aversion.

We can recreate our responses to the world around us into an experience that's a lot more pleasant to be with. It's still a constructed reality and we'll need to let it unravel completely at some point in the process of deconstructing our way of being in the world. This will lead to a much more sublime, insightful, and penetrating experience leading to the ending of dukkha. The mind that has been repatterned into a more pleasant, open, expansive state of mind allows us to experience the cessation of dukkha much more easily.

Over the years, I've used contemplation as a way of bringing up skillful perceptions of the opposites of the five hindrances. I pose a question to myself, such as, ``What would it feel like if I were free from sensual desire and I didn't want anything?'' One word that comes up is ``contentment.'' What's it like? We've all had moments of contentment when we don't want anything, when desire isn't arising and we're at peace. The mind isn't surging out towards picking up sensory experience in the realm of sights, sounds, tastes, or tactile sensations. The mind is at rest and absolutely happy with what is happening right now. Just imagine what that's like, even for a moment. Imagine the feeling of contentment.

\enlargethispage{1\baselineskip}
What is it like to be free from ill will or aversion, when we have let go of something that we were gnawing on for a long time? This could be a memory of a conversation or a constructed perception of a person who has given us a hard time. We realize the dukkha in that and let it go or establish an alternative perception. We can create an image that brings a feeling of kindness, gentleness, or open-heartedness. We can tag it with a label like ``\emph{mettā} (loving-kindness)'' or ``gentle.'' What's a good word to describe that quality for yourself?

One enjoyable memory for me is again about my mother. When I was going through some of the photos of my \linebreak mother, there was one of her at the nursing home, an image that is very much embedded in my mind. She was experiencing a moderate amount of memory loss and some delirium and was sitting across a table from me. There were a few other people around, including my sister and one of the nurses. My mother was kind of out of it, and we were all talking and engaging with each other. At one point I tried to include her by asking her how she was doing. I wasn't sure whether she was mentally present or not. She didn't look very alert and didn't really engage. I caught her eye and said, ``Mom, are you following all of this? Are you following any of this?'' She just looked a little confused. And then I gently said, ``Do you know who I am?'' She responded with another momentary look of confusion, but then all of a sudden she became alert and was 100 percent present with us. Her big, blue eyes stared out from her wrinkled, ninety-two-year-old face, and with a penetrating gaze, directly into my eyes, she said slowly, ``You are mine.'' And my heart just melted.

In a moment of aversion or going through a rough patch, if I can bring that memory up, it immediately washes that mood away. Instead of aversion, I feel kindness and gentleness. The world opens up again.

We can change our perceptions on the spot by having a word or an image that we can pull up. We can ask ourselves, ``What would it be like to be free from ill will? What would it be like to be free from dullness and drowsiness? What is it like to have a fresh, alert mind?'' You can think of the words ``alert,'' ``fresh.'' What is the feeling that arises? Cool and spacious.{\parfillskip=0pt \par}

If you're feeling agitated, ask yourself, ``What's the opposite like? What's it like when you're not agitated? Still, silent, quiet?'' We've all had moments of peace, even if they've been very fleeting. What's the feeling that the word ``stillness'' brings into perceptual experience?

What is it like to be free from doubt? We feel secure, grounded, stable, unwavering. What does that feel like in the body and in the mind? This self-questioning, then labeling and experiencing the feeling of the antidote, are ways that we reconstruct these perceptual patterns.

When we identify which hindrances are present and experience them directly, we can bring up perceptions of their opposites as a meditation. You can do that as a mantra. After you've established yourself comfortably in the sitting posture, take a few moments to say, ``Okay, what's it like to be contented, gentle, alert? Repeat the words: ``silent,'' ``still,'' ``secure.'' Watch what happens to the heart. It opens, relaxes, and settles. As that happens, the opportunity arises to be present with something simple and soothing, like the breathing. If the mind is not agitated from one of the hindrances, but instead is experiencing positive qualities, this is the mind that composes itself easily.

Freedom from the hindrances can come through reconfiguring our perceptual patterns. It's also important to remind ourselves that these patterns aren't who we are, even the most positive ones. We don't need to take ownership of them and say, ``Finally, I've gotten rid of my self that was full of hindrances, and now I am someone who is full of more positive qualities.'' We're not trying to construct an identity around them either. We're trying to temporarily fabricate skillful states that lead us to a deeper understanding. We create the mind space that can understand the arising and passing of the world around us and how to deconstruct all of our perceptual patterns.

We have to start by recognizing the mind states that obstruct us. We can repattern them in a way that causes us to open and relax and make life enjoyable. We have to enjoy life. Even though we're working towards deconstructing how we view the world and ourselves within it, we have to have patterns that are pleasant to experience. We can then live with ourselves more easily.

We can be an example for others, too. We become easier to be around if we work through our obstructive patterns and create skillful ones that appear to the world around us. It gives other people hope. We naturally want to live in the world with ease of mind, and this will also help other people.

The most obstructive perception that we live with is the perception of self. We can't really help other people fully until we have let go of all selfing patterns. We've got to get out of the way if we really want to be of service to the world. If we can learn how to create a deep sense of contentment and happiness in ourselves through eliminating obstructive patterns, particularly the self-perception, then we can be of greater benefit to the world. That's my aspiration anyway. The potential to make these changes is what keeps me here after twenty-five plus years. There is still lots of work to do, but many benefits have come along the way. It gives me hope, faith, and inspiration for the future.
